% Part: first-order-logic
% Chapter: proof-systems
% Section: introduction

\documentclass[../../../include/open-logic-section]{subfiles}

\begin{document}

\iftag{FOL}
      {\olfileid{fol}{prf}{int}}
      {\olfileid{pl}{prf}{int}}

\olsection{Pambuka}

Logika lumrahe duwe semantik lan sistem !!a{derivation}. Semantik
ngrembug konsep kayata kabeneran, bisa dicukupi, validitas, lan akibat
semantis. Tujuwane sistem !!{derivation} yaiku menehi metodhe kang
murni sintaktis kanggo netepake akibat semantis lan validitas. Sistem
kasebut murni sintaktis ing teges yen !!{derivation} ing sistem kaya
mangkono iku obyek sintaktis winates, lumrahe urutan (utawa susunan
winates liyane) saka !!{sentence}s utawa !!{formula}s. Sistem
!!{derivation} kang becik nduweni sipat yen saben urutan utawa susunan
!!{sentence}s utawa !!{formula}s kang diwenehake bisa dipriksa kanthi
mekanis apa ``bener''.

Sistem !!{derivation} kang paling prasaja (lan paling dhisik saka sisi
sejarah) kanggo logika tataran kapisan iku \emph{aksiomatik}. Urutan
!!{formula}s dianggep minangka !!a{derivation} ing sistem kaya mangkono
yen saben !!{formula} ing jerone kalebu sawijining himpunan ``aksioma''
kang tetep, utawa ndherek saka !!{formula}s sadurunge ing urutan
nganggo salah siji saka sawetara ``aturan inferensi'' kang tetep---lan
bisa dipriksa kanthi mekanis apa !!a{formula} iku aksioma lan apa
formula kasebut kanthi bener ndherek !!{formula}s liyane nganggo salah siji
aturan inferensi. Sistem !!{derivation} aksiomatik gampang dijlentrehake---
lan uga gampang ditangani saka sisi metateori---nanging !!{derivation}s
ing jerone angel diwaca lan dimangerteni, uga angel digawe.

Sistem !!{derivation} liyane wis dikembangake supaya luwih gampang
mbangun !!{derivation}s utawa luwih gampang mangerteni !!{derivation}s
sawise rampung. Tuladhane yaiku deduksi natural, wit kabeneran kang uga
diarani bukti tableau, lan kalkulus sekuen. Sawetara sistem
!!{derivation} dirancang mligi kanthi nggatekake mekanisasi; upamane,
metodhe resolusi gampang ditindakake ing piranti alus (nanging
!!{derivation}s-e pokoke mokal dimangerteni). Umume sistem
!!{derivation} liyane iki nggambarake !!{derivation}s minangka wit
!!{formula}s lan dudu urutan. Kanthi mangkono, luwih gampang katon
perangan !!a{derivation} endi gumantung marang perangan liyane kang endi.

Mula, kanggo sawijining logika, kayata logika tataran kapisan, sistem
!!{derivation} kang beda-beda bakal menehi jlentrehan kang beda ngenani
apa tegese !!a{sentence} dadi \emph{teorema} lan apa tegese
!!a{sentence} bisa !!{derivable} saka sawetara ukara liyane. Apa wae
carane (liwat !!{derivation}s aksiomatik, deduksi natural,
!!{derivation}s sekuen, wit kabeneran, utawa refutasi resolusi), kita
ngarepake relasi kasebut cocog karo gagasan semantis validitas lan
akibat semantis. Kita nulis $\Proves !A$ kanggo ``$!A$~iku teorema''
lan ``$\Gamma \Proves !A$'' kanggo ``$!A$~bisa !!{derivable}
saka~$\Gamma$.'' Kepriye wae $\Proves$~ditemtokake, kita ngarepake
simbol iku cocog karo $\Entails$, yaiku:
\begin{enumerate}
\item $\Proves !A$ yen lan mung yen $\Entails !A$
\item $\Gamma \Proves !A$ yen lan mung yen $\Gamma \Entails !A$
\end{enumerate}
Arah ``mung yen'' ing ndhuwur diarani \emph{kasahihan}.
!!^a{derivation} sistem iku sah yen !!{derivability} njamin akibat
semantis (utawa validitas). Saben sistem !!{derivation} kang patut
kudu sah; sistem !!{derivation} kang ora sah ora migunani babar pisan.
Awit kabeh tujuwane !!a{derivation} yaiku menehi jaminan sintaktis
tumrap validitas utawa akibat semantis. Kita bakal mbuktekake kasahihan
kanggo sistem !!{derivation} kang disajekake.

Arah walikan, yaiku ``yen'', uga wigati: arah iki diarani
\emph{kelengkapan}. Sistem !!{derivation} kang jangkep cukup kuwat
kanggo nuduhake yen $!A$~iku teorema saben $!A$~valid, lan yen
$\Gamma \Proves !A$ saben $\Gamma \Entails !A$. Kelengkapan luwih angel
ditetepake, lan sawetara logika ora duwe sistem !!{derivation} kang
jangkep. Logika tataran kapisan duwe. Kurt G\"odel dadi wong kapisan
kang mbuktekake kelengkapan kanggo !!a{derivation} sistem logika
tataran kapisan ing disertasine taun 1929.

Konsep liyane kang gegayutan karo sistem !!{derivation} yaiku
\emph{konsistensi}. Himpunan !!{sentence}s diarani ora konsisten yen
apa wae bisa !!{derive}d saka himpunan iku, lan diarani konsisten yen
ora mangkono. Ora konsisten iku pasangan sintaktis saka ora bisa
dicukupi: kaya himpunan kang ora bisa dicukupi, himpunan !!{sentence}s
kang ora konsisten ora dadi teori kang becik, amarga cacat kanthi
dhasar. Himpunan !!{sentence}s kang konsisten durung mesthi bener utawa
migunani, nanging saora-orane wis ngliwati ambang paling cilik kanggo
piguna logis. Kanggo sistem !!{derivation} kang beda, definisi mligi
ngenani konsistensi himpunan !!{sentence}s bisa beda, nanging kaya
$\Proves$, kita ngarepake konsistensi cocog karo pasangan semantise,
yaiku bisa dicukupi. Kita ngarepake tansah lumaku yen $\Gamma$
konsisten yen lan mung yen bisa dicukupi. Ing kene, arah ``mung yen''
padha karo kelengkapan (konsistensi njamin bisa dicukupi), lan arah
``yen'' padha karo kasahihan (bisa dicukupi njamin konsistensi).
Nyatane, kanggo logika klasik tataran kapisan, rong versi kasahihan lan
kelengkapan iki padha ekuivalen.

\end{document}
